\documentclass[aspectratio=169,11pt]{beamer}
\usepackage{teaching_slides}
\usepackage{listings}

\lstset{
  language=R,
  basicstyle=\ttfamily\tiny,
  backgroundcolor=\color{gray!5},
  frame=single,
  rulecolor=\color{gray!50},
  keywordstyle=\color{blue!70!black},
  commentstyle=\color{green!40!black},
  stringstyle=\color{orange!90!black},
  columns=flexible,
  keepspaces=true,
  showstringspaces=false,
  breaklines=true
}


\usepackage{inconsolata} % nice monospace font


\title{Program Evaluation and Randomized Experiments}
\author{Chris Conlon}
\institute{Applied Econometrics}
\date{\today}

\begin{document}


\frame{\titlepage}

\frame{\frametitle{Overview}
This set of lectures will cover (roughly) the following papers:\\
Theory:
\begin{itemize}
\item Angrist and Imbens (1994)
\item Heckman Vytlacil (2005/2007)
\item Abadie and Imbens (2006)
\end{itemize}
And draw heavily upon notes by 
\begin{itemize}
\item Guido Imbens
\item Richard Blundell and Costas Meghir
\end{itemize}

}

\begin{frame}
\frametitle{The Evaluation Problem}
\begin{itemize}
\item The issue we are concerned about is identifying the effect of a policy or an investment or some individual action on one or more outcomes of interest
\item This has become the workhorse approach of the applied microeconomics fields (Public, Labor, etc.)
\item Examples may include:
\begin{itemize}
\item The effect of taxes on labor supply
\item The effect of education on wages
\item The effect of incarceration on recidivism
\item The effect of competition between schools on schooling quality
\item The effect of price cap regulation on consumer welfare
\item The effect of indirect taxes on demand
\item The effects of environmental regulation on incomes
\item The effects of labor market regulation and minimum wages on wages and employment
\end{itemize}
\end{itemize}
\end{frame}

\begin{frame}{Potential Outcomes}
\begin{itemize}
\item Consider a binary treatment $D_i \in \{0,1\}$. 
\begin{itemize}
\item Some people use $T_i \in \{0,1\}$ instead.
\end{itemize}
\item We observe the outcome $Y_i$. But there are two \alert{potential outcomes}
\begin{itemize}
\item $Y_i(1)$ the outcome for $i$ if they  \alert{are treated}.
\item $Y_i(0)$ the outcome for $i$ if they \alert{are not treated} (control).
\end{itemize}
\item We are generally interested in $\beta_i \equiv  Y_i(1)-Y_i(0)$ which we call the \alert{treatment effect}.
\begin{itemize}
\item Individuals have \alert{heterogeneous treatment effects}.
\item In an ideal world we could fully characterize $f(\beta_i)$
\end{itemize}
\end{itemize}
\end{frame}

\begin{frame}{Some Challenges}
Stable unit treatment value assumption (SUTVA)
\begin{itemize}
\item We assume a \textit{ceteris paribus} version of treatment effects
\item We need $\beta_i$ to be a policy invariant (structural) parameter.
\item Your $\beta_i$ doesn't respond to whether or not another individual is treated.
\item Two common limitations:
\begin{itemize}
\item Peer effects: Whether you respond to job training program depends on whether your spouse is also treated.
\item Equilibrium effects: if we sent everyone to college, returns to college would be quite different.
\end{itemize}
\end{itemize}
\end{frame}



\begin{frame}{Some Challenges}
\begin{columns}[T] % align columns
\begin{column}{.65\textwidth}
Fundamental Problem of Causal Inference
\begin{itemize}
\item We don't observe the \alert{counterfactual} $Y_i(D_i)$.
\item For a single individual we either observe $Y_i(1)$ \alert{or} $Y_i(0)$ but never both!
\begin{itemize}
\item ex: We don't see what your wage would have been if you didn't attend college.
\item ex: We might know your cholesterol before you took Lipitor, but we don't know what it would be today if you didn't take Lipitor.
\end{itemize}
\end{itemize}
\end{column}%
\hfill%
\begin{column}{.38\textwidth}

  \vspace{20pt}
  
  $Y_{i} = D_{i}\cdot Y_{i}(1) + (1-D_{i}) \cdot Y_{i}(0)$\\
  \vspace{20pt}
  \begin{tabular}{ccccc}
    \toprule
    i & $Y_{i}(1)$ &  $Y_{i}(0)$ & $D_{i}$ & $D_{i}$ \\
    \midrule
    1 &     1      &     \alert{?}      &   1   & 1 \\
    2 &     0      &     \alert{?}      &   1   & 0 \\
    3 &     \alert{?}      &     0      &   0   & 0 \\
     & &  $\vdots$ & & \\
    $n$ &     \alert{?}      &     1      &   0   & 1 \\    
  \end{tabular}
\end{column}%
\end{columns}
\end{frame}


\end{document}